\section{An Intervention Audit of Reusable Procedures}
\label{sec:mechanism}
The audit treats a reusable procedure as an intervention rather than a label attached to a successful run. Let endpoint $i$ specify a question $x_i$, evidence package $E_i$, model configuration $M_i$, and family-specific binary outcome $Y_i\in\{0,1\}$. For repeated executions under arm $a$, let $\bar Y_i(a)$ be the repeat mean. Repeats measure execution robustness; endpoints remain the independent scientific units.

\paragraph{Base stress test: N, G, and T.}
The historical design contains the original agent N, a complexity-matched target-blind \emph{stress helper} G, and the targeted procedure T. G is intentionally mechanism-misaligned rather than a benign or optimal generic assistant; its role is to test whether T-versus-structured-assistance can produce a misleading verdict. We report
\begin{equation}
\widehat{\Delta}_{T-N}=\frac{1}{|\mathcal I|}\sum_i\bigl(\bar Y_i(T)-\bar Y_i(N)\bigr),\qquad
\widehat{\Delta}_{T-G}=\frac{1}{|\mathcal I|}\sum_i\bigl(\bar Y_i(T)-\bar Y_i(G)\bigr).
\end{equation}
T--N is the net effect of the targeted intervention relative to the original agent. T--G asks whether it beats the executed target-blind helper. G is not assumed benign, so T$>$G with T$\leq$N is no repair.

\paragraph{Same-surface placebo: G$_0$.}
G$_0$ executes the same helper-exposure path as T but is information-empty: it reads neither package nor context and returns an empty object. We report
\begin{equation}
\widehat{\Delta}_{T-G_0}=\frac{1}{|\mathcal I|}\sum_i\bigl(\bar Y_i(T)-\bar Y_i(G_0)\bigr),\qquad
\widehat{\Delta}_{G_0-N}=\frac{1}{|\mathcal I|}\sum_i\bigl(\bar Y_i(G_0)-\bar Y_i(N)\bigr).
\end{equation}
T--G$_0$ is a \emph{same-surface operation-output contrast}: it asks whether supplying the targeted output changes behavior relative to no operation output under the same exposure surface. G$_0$--N is the corresponding surface/placebo perturbation diagnostic. These quantities obey the observed-mean identity
\begin{equation}
\widehat{\Delta}_{T-N}=\widehat{\Delta}_{T-G_0}+\widehat{\Delta}_{G_0-N}.
\label{eq:decomp}
\end{equation}
Equation~\ref{eq:decomp} is a decomposition of contrasts, not a mediation theorem. When G$_0$--N is nonzero, T--G$_0$ is conditional on a perturbed surface and must not be called a pure operation effect. Accordingly, we classify T--N$>0$ as \emph{net repair}; T--G$_0$ is reported separately as evidence about the output under that surface. Cells with nonzero placebo shift are explicitly labeled surface-interacting rather than operation-attributed.

\paragraph{Integration-surface control: R$_{\mathrm{surf}}$.}
R$_{\mathrm{surf}}$ runs the exact same frozen operation source on the exact same candidate evidence as T, then materializes the identical operation output inside ordinary context while exposing no callable helper. Thus
\begin{equation}
\widehat{\Delta}_{T-R_{\mathrm{surf}}}=\frac{1}{|\mathcal I|}\sum_i\bigl(\bar Y_i(T)-\bar Y_i(R_{\mathrm{surf}})\bigr)
\end{equation}
is an integration-surface estimand under the frozen forced one-answer harness. R$_{\mathrm{surf}}$ is deliberately not an independently learned or heuristic temporal retriever. A zero-call parity audit verifies that candidate evidence is preserved, operation-output content is identical, and only where the precomputed output enters the harness changes. The prospectively frozen material threshold is 10 percentage points on the pre-frozen DeepSeek endpoint portfolio. Multi-turn deployments in which a callable helper can be invoked adaptively require a different control.

\begin{figure}[t]
\centering
\includegraphics[width=\linewidth]{figures/fig1_temporal_skill_causal.pdf}
\caption{Base three-arm stress test. N anchors the original agent, G is a target-blind mechanism-misaligned stress helper, and T contains the targeted temporal operation. G$_0$ exposes placebo/surface perturbation; R$_{\mathrm{surf}}$ isolates where the same precomputed operation output enters the forced one-answer harness.}
\label{fig:causal}
\end{figure}

The audit targets three recurring failure families. \textbf{Cutoff validity} filters evidence by the target information cutoff. \textbf{Release alignment} maps a release to the reporting period represented by its content. \textbf{Exogenous grounding} links a numerical change to first-party documentary evidence available before the cutoff. No targeted procedure contains endpoint-specific answers, gold labels, or fitted values; procedure hashes are frozen before execution and reused on held-out releases.

Three falsifiers follow. First, T$>$G with T$\leq$N is no repair. Second, a positive T--G$_0$ does not by itself establish a pure operation effect when G$_0$--N is nonzero; all three contrasts in Eq.~\ref{eq:decomp} must remain visible. Third, if the integration surface matters under the forced one-answer harness, T should leave a material residual over R$_{\mathrm{surf}}$ when operation output and evidence are held exactly fixed.
